% !TEX root = chapter-standalone.tex

\section{Twisted interval constructions}
\linkvideo{spring2021-tradeoffs:tradeoffs:orders:interval-poset} % Poset of intervals


\todotext{2026-07 JL: Explain that there are TWO twisted interval constructions, and that they are opposites of each other.}

\begin{marginfigure}
    \centering
    %\includesag{first_poset_intervals}
    \includesag{tw_int_fig}
    \caption{
        $\interv \posela \poselb
            \posleqof{\posint{\posA}}
            \interv \poselc \poseld$
    }
\end{marginfigure}

\todotext{
    We should define there $ \udpL$ and $ \udpU$.
}

\begin{definition}[``Twisted'' preorder of intervals~$\posint\posA$]
    \label{def:twisted-poset-of-intervals}
    \SYNDEF{twisted preorder of intervals}
    Given a preorder~\posA, we define a \emph{``twisted'' \SY{preorder} of intervals}~$\posint\posA$, whose elements are the intervals of~\posA,
    by ordering the intervals by inclusion of their associated subsets:
    %
    \begin{equation}\label{eq:twisted-1}
        \prfdoubleperiod{
            \interv \posela \poselb
            \posleqof{\posint{\posA}}
            \interv \poselc \poseld
        }{
            \interv \posela \poselb
            \setsubseteq
            \interv \poselc \poseld
        }
    \end{equation}
    Equivalently we only need to check the bounds:
    \begin{equation}\label{eq:twisted-2}
        \prfdoubleperiod{
            \interv \posela \poselb
            \posleqof{\posint{\posA}}
            \interv \poselc \poseld
        }{
            (\poselc \posAleq \posela ) \booland (\poselb \posAleq \poseld)
        }
    \end{equation}

\end{definition}

\begin{exercise}
    Check that the relation defined in \cref{def:twisted-poset-of-intervals} is reflexive and transitive, so that~$\posint\posA$ is indeed a \SY{preorder}.
    Then show that if~\posA is a \SY{poset}, then~$\posint\posA$ is a \SY{poset}.
\end{exercise}
\begin{solution}
    Write~$\posleq$ for~$\posleqof{\posint{\posA}}$, and use the characterization~\eqref{eq:twisted-2}.
    \begin{itemize}
        \item Reflexivity: for any interval~$\interv{\posela_1}{\poselb_1}$, we have~$\posela_1\posAleq \posela_1$ and~$\poselb_1\posAleq \poselb_1$ by reflexivity of~$\posAleq$, and therefore~$\interv{\posela_1}{\poselb_1} \posleq \interv{\posela_1}{\poselb_1}$.
        \item Transitivity: suppose~$\interv{\posela_1}{\poselb_1} \posleq \interv{\posela_2}{\poselb_2}$ and~$\interv{\posela_2}{\poselb_2} \posleq \interv{\posela_3}{\poselb_3}$.
              By definition, this means~$\posela_2 \posAleq \posela_1$,~$\posela_3 \posAleq \posela_2$,~$\poselb_1 \posAleq \poselb_2$, and~$\poselb_2 \posAleq \poselb_3$.
              By transitivity of~$\posAleq$, we get~$\posela_3 \posAleq \posela_1$ and~$\poselb_1 \posAleq \poselb_3$, that is,
              \begin{equation}
                  \interv{\posela_1}{\poselb_1} \posleq \interv{\posela_3}{\poselb_3}.
              \end{equation}
        \item Antisymmetry (when~\posA is a \SY{poset}): suppose~$\interv{\posela_1}{\poselb_1} \posleq \interv{\posela_2}{\poselb_2}$ and~$\interv{\posela_2}{\poselb_2} \posleq \interv{\posela_1}{\poselb_1}$.
              Then~$\posela_2 \posAleq \posela_1$ and~$\posela_1 \posAleq \posela_2$, so~$\posela_1=\posela_2$ by antisymmetry of~$\posAleq$; in the same way,~$\poselb_1=\poselb_2$.
              Hence~$\interv{\posela_1}{\poselb_1} = \interv{\posela_2}{\poselb_2}$.
    \end{itemize}
\end{solution}

In general, $\posint\posA$ does not have a top or a bottom.

\vfill
\clearpage

\subsection{Set-based filtering}
\linkvideo{spring2021-tradeoffs:tradeoffs:orders:set-based-filtering} % Set-based filtering
\todojira{458}{\bernina:Expand.
    Need cute pictures of dynamics.
}
We now look at an example of \textbf{set-based filtering}, where filtering refers to online inference (recursive estimation).
Suppose that we want to track the value of a quantity~$x\setin \interv{0}{100}$, without having \emph{a priori} information about~$x$.
We are equipped with sensors, which periodically measure the quantity~$x$ with some variable precision.
At time~$t\setin \nonNegReals$ they produce an \emph{observation}~$y_t\colon x_t\setin \interv{l_t}{u_t}$.
Also, note that the quantity fluctuates randomly, and we bound its ``velocity'' to be~$\dot{x}_t\setin \interv\minusone\plusone$ (except at boundaries).
At the beginning, our information state~$\bar{i}_0$ could be that~$x\setin \interv{0}{100}$.
At time 0, we get an observation~$y_0$, that says~$x\setin \interv{21}{24}$.
The new information state can be obtained by ``fusing'' the two inputs we have received about~$x$.
This corresponds to the intersection
\begin{equation}\label{eq:filtering-2}
    \prfdoubleperiod{
        x\setin \pars{ \interv{0}{100} \setintersection \interv{21}{24}}
    }{
        x\setin \interv{21}{24}
    }
\end{equation}
Say we get an observation~$y_1$ which says~$x\setin \interv{19}{22}$.
\begin{figure}
    \centering
    \includesag{080_hasse_filtering_bis}
    \caption{}
    \label{fig:hasse_filtering}
\end{figure}
We now need to take into account the evolution/dynamics of the quantity we are tracking.
From the interval~$\interv{21}{24}$ we know that the variable could have evolved in~$\interv{20}{25}$ (dynamics are bounded with a unit increase/decrease).
Therefore, the new information state is given by
\begin{equation}\label{eq:filtering-3}
    \prfdoubleperiod{
        x\setin \pars{ \interv{20}{25} \setintersection \interv{19}{22}}
    }{
        x\setin \interv{20}{22}
    }
\end{equation}
One of the structures which could sustain this kind of inference, is the twisted \SY{preorder} of intervals (\cref{def:twisted-poset-of-intervals}).

The Hasse diagram representing a situation related to this example could be as reported in \cref{fig:hasse_filtering}.

